% Ad Familiares 11.20 (ad-familiares-11-20)
% Source: https://raw.githubusercontent.com/PerseusDL/canonical-latinLit/master/data/phi0474/phi056/phi0474.phi056.perseus-lat1.xml
% Fetched by scripts/fetch_latin.py

% Dateline (Perseus): Scr. Eporediae ix K. Iun. a. 711 (43).

\ciceroLetterOpener{D. BRVTVS S. D. M. CICERONI}

\ciceroSection{1} quod pro me non facio, id pro te facere amor meus in te tuaque officia cogunt, ut timeam. saepe enim mihi cum esset dictum neque a me contemptum, novissime Labeo Segulius, homo sui simillimus, narrat mihi apud Caesarem se fuisse multumque sermonem de te habitum esse. ipsum Caesarem nihil sane de te questum, nisi dictum quod diceret te dixisse 'laudandum adulescentem, ornandum, tollendum' ; se non esse commissurum ut tolli possit. hoc ego Labeonem credo illi rettulisse aut finxisse dictum, non ab adulescente prolatum. veteranos vero pessime loqui volebat Labeo me credere et tibi ab iis instare periculum maximeque indignari, quod in decem viris neque Caesar neque ego habiti essemus atque omnia ad vestrum arbitrium essent conlata.

\ciceroSection{2} haec cum audissem et iam iri itinere essem, committendum non putavi prius ut Alpis transgrederer quam quid istic ageretur scirem ; nam de tuo periculo crede mihi iactatione verborum et denuntiatione periculi sperare eos te pertimefacto, adulescente impulso, posse magna consequi praemia, et totam istam cantilenam ex hoc pendere ut quam plurimum lucri faciant. neque tamen non te cautum esse volo et insidias vitantem ; nihil enim tua mihi vita potest esse iucundius neque carius.

\ciceroSection{3} illud vide ne timendo magis timere cogare et quibus rebus a potest occurri veteranis occurras, primum quod desiderant de decem viris facias, deinde de praemiis, si tibi videtur, agros eorum militum qui cum Antonio veterani fuerunt iis dandos censeas ab utrisque nobis ; de nummis lente ac ratione habita pecuniae senatum de ea re constituturum. quattuor legionibus iis, quibus agros dandos censuistis, video facultatem fore ex agris Sullanis et agro Campano ; aequaliter aut sorte agras legionibus adsignari puto oportere.

\ciceroSection{4} haec me tibi scribere non prudentia mea hortatur sed amor in te et cupiditas oti, quod sine te consistere non potest. ego , nisi valde necesse fuerit, ex Italia non excedam. legiones armo, paro. spero me non pessimum exercitum habiturum ad omnis casus et impetus hominum. de exercitu quem Pansa habuit legionem mihi Caesar non remittit. ad has litteras statim mihi rescribe tuorumque aliquem mitte, si quid reconditum magis erit meque scire opus esse putaris. viiii K. Iun. Eporedia.
